\chapter{Real/первопорядковый допуск минимального \texorpdfstring{$R_2$}{R2}-коннектора}
\label{ch:v7-r2-real-first-order-admission-gate}

\section{Почему графовой минимальности недостаточно}

Предыдущая глава классифицировала дополнения абстрактного двудольного графа
хиральностей и выделила единственный одно-мультиплетный кандидат
$R_2\sim(\mathbf3,\mathbf2)_{7/6}$. Теперь требуется вернуть полные
бимодульные метки алгебры Стандартной модели
\begin{equation}
 \mathcal A_{\rm SM}=\mathbb C\oplus\mathbb H\oplus M_3(\mathbb C).
 \label{eq:v7-r2-standard-algebra}
\end{equation}
Граф по хиральности забывает, какая компонента алгебры действует слева и
справа. Именно эти две координаты проверяет условие первого порядка.

\section{Точный критерий ребра}

Для неприводимого бимодуля $\mathcal H_{ij}$ и блока
$D_{ij,kl}:\mathcal H_{ij}\to\mathcal H_{kl}$ двойной коммутатор содержит
\begin{equation}
 [[D,a],JbJ^{-1}]_{ij,kl}
 =(a_i-a_k)D_{ij,kl}(b_j-b_l).
 \label{eq:v7-r2-first-order-factor}
\end{equation}
Поэтому для всех $a,b$ ненулевое ребро допускается только при
\begin{equation}
 i=k\qquad\text{или}\qquad j=l.
 \label{eq:v7-r2-same-row-column-criterion}
\end{equation}
Достаточно проверить центральные идемпотенты трёх слагаемых алгебры: если
меняются обе координаты, можно независимо выбрать ненулевыми оба множителя
в \eqref{eq:v7-r2-first-order-factor}.

\section{Бимодульные координаты заряженного пакета}

На уровне слагаемых алгебры стандартные вершины имеют координаты
\begin{align}
 Q_L&=(\mathbb H,M_3),& L_L&=(\mathbb H,\mathbb C),\nonumber\\
 u_R&=(\mathbb C,M_3),& d_R&=(\mathbb C,M_3),&
 e_R&=(\mathbb C,\mathbb C).
 \label{eq:v7-r2-bimodule-coordinates}
\end{align}
Различие обычного и сопряжённого действия $\mathbb C$ не влияет на
проверяемое препятствие по слагаемым.

Три существующих юкавских ребра меняют ровно одну координату:
\begin{align}
 Q_L\to u_R&:(\mathbb H,M_3)\to(\mathbb C,M_3),\nonumber\\
 Q_L\to d_R&:(\mathbb H,M_3)\to(\mathbb C,M_3),\nonumber\\
 L_L\to e_R&:(\mathbb H,\mathbb C)\to(\mathbb C,\mathbb C).
 \label{eq:v7-r2-existing-edge-pass}
\end{align}
Они проходят критерий \eqref{eq:v7-r2-same-row-column-criterion}.

Все отсутствующие рёбра меняют обе координаты:
\begin{align}
 L_L\to u_R&:(\mathbb H,\mathbb C)\to(\mathbb C,M_3),\nonumber\\
 L_L\to d_R&:(\mathbb H,\mathbb C)\to(\mathbb C,M_3),\nonumber\\
 Q_L\to e_R&:(\mathbb H,M_3)\to(\mathbb C,\mathbb C).
 \label{eq:v7-r2-missing-edge-fail}
\end{align}
Следовательно, не только выбранная $R_2$-пара, но и каждое минимальное
прямоугольное дополнение предыдущей главы содержит запрещённые диагонали.

\section{Почему Real-дополнение не исправляет дефект}

Реальная структура переводит
\begin{equation}
 J:\mathcal H_{ij}\longmapsto\mathcal H_{ji}.
 \label{eq:v7-r2-real-coordinate-swap}
\end{equation}
Если исходное ребро меняет $i\ne k$ и $j\ne l$, сопряжённое ребро меняет
$j\ne l$ и $i\ne k$. Поэтому добавление $JD_{R_2}J^{-1}$ восстанавливает
самосопряжённость и Real-парность, но не обнуляет двойной коммутатор.

\section{Граница \texorpdfstring{$S^0$}{S0}-реальности и литературы}

В стандартной конечной геометрии строгая $S^0$-реальность сохраняет
разделение частиц и античастиц и исключает прямой кварк-лептонный сектор.
После отказа от неё в ранней модели возникал иной канал
$L_L\leftrightarrow u_R^c$, но не двухрёберный прямоугольник текущего
$R_2$-кандидата; полученная модель также столкнулась с феноменологическими
ограничениями \cite{v7r2-paschke}. Современные расширенные скалярные сектора
возникают после явного изменения геометрической архитектуры, например
Clifford-условий или обобщённых флуктуаций
\cite{v7r2-kurkov,v7r2-generalized}.

Это означает, что проверка спектрального действия и цветового вакуума для
неизменённого родителя не начинается: требуемого допустимого блока $D_F$ в
нём нет.

\begin{theorem}[Первопорядковый запрет фиксированного $R_2$]
В стандартной алгебре $\mathbb C\oplus\mathbb H\oplus M_3(\mathbb C)$ на
фиксированных вершинах $H_{15}$ каждое отсутствующее кварк-лептонное ребро
меняет одновременно левую и правую бимодульные координаты. Поэтому оба
ребра графово минимального $R_2$-дополнения и их Real-сопряжения нарушают
строгое условие первого порядка.
\end{theorem}

\begin{nogocriterion}[Строгий неизменённый родитель закрыт]
Нельзя переходить к потенциалу $R_2$ или цветному вакууму, сохраняя
одновременно стандартную алгебру и представление, фиксированные вершины,
обычную Real-структуру и строгое условие первого порядка. Сначала необходимо
явно изменить хотя бы один из этих пунктов.
\end{nogocriterion}

\begin{protocol}
\begin{enumerate}
 \item существующие юкавские рёбра проходят первый порядок: \textbf{да};
 \item $L_L-u_R$ проходит: \textbf{нет};
 \item $Q_L-e_R$ проходит: \textbf{нет};
 \item Real-сопряжения устраняют дефект: \textbf{нет};
 \item какое-либо отсутствующее ребро допустимо на фиксированных вершинах:
 \textbf{нет};
 \item строгий $R_2$-родитель допущен: \textbf{нет};
 \item стадия спектрального действия и цветового вакуума достигнута:
 \textbf{нет};
 \item следующий гейт: \textbf{минимальное расширение алгебры,
 представления или скрученного/обобщённого первого порядка}.
\end{enumerate}
\end{protocol}

Машинный сертификат сохранён в
\path{s2t_v7_r2_real_first_order_admission_gate_results.json}.

\begin{thebibliography}{9}
\bibitem{v7r2-paschke}
M.~Paschke, F.~Scheck, A.~Sitarz,
\emph{Can (noncommutative) geometry accommodate leptoquarks?},
Physical Review D 59 (1999), 035003; arXiv:hep-th/9709009.
\bibitem{v7r2-kurkov}
M.~A.~Kurkov, F.~Lizzi,
\emph{Clifford Structures in Noncommutative Geometry and the Extended Scalar Sector},
Physical Review D 97 (2018), 085024; arXiv:1801.00260.
\bibitem{v7r2-generalized}
A.~H.~Chamseddine, A.~Connes, W.~D.~van Suijlekom,
\emph{Inner Fluctuations in Noncommutative Geometry without the First Order Condition},
Journal of Geometry and Physics 73 (2013), 222--234; arXiv:1304.7583.
\end{thebibliography}